\begin{frame}{Faithfulness: perturb important features}
\HierarchyHeader{\ContentIcon}{Content}{Correctness / faithfulness}{Do important attributed features control the model output?}
Let $a=E(f,x)$, where $a_i$ is feature importance. A perturbation-based test asks whether replacing features with large $|a_i|$ changes $f(x)$.
\[
\Delta_i=f(x)-f(\xdel{i})
\]
Here $\xdel{i}$ denotes $x$ with feature $i$ replaced by the chosen baseline. The replacement rule is part of the metric.
\end{frame}

\begin{frame}{Faithfulness correlation}
\HierarchyHeader{\ContentIcon}{Content}{Correctness / faithfulness}{Does attribution mass correlate with output change?}
\[
\operatorname{FC}(x)=\operatorname{corr}_{S}\left(\sum_{i\in S}a_i,\; f(x)-f(\xdel{S})\right)
\]
Positive association is evidence that attribution mass tracks output change under the sampled subsets and replacement protocol.
\end{frame}

\begin{frame}{Deletion curve}
\HierarchyHeader{\ContentIcon}{Content}{Correctness / faithfulness}{What happens when top-ranked features are removed first?}
Rank features by $|a_i|$, then progressively replace the most important ones. A steep score decrease early in the curve is consistent with a useful ranking under this deletion rule.
\centering\includegraphics[width=.7\linewidth,height=.46\textheight,keepaspectratio]{deletion_curve.pdf}
\end{frame}

\begin{frame}{Deletion AUC and AOPC}
\HierarchyHeader{\ContentIcon}{Content}{Correctness / faithfulness}{How do we summarize a deletion trajectory?}
\[
\mathrm{AUC}_{del}=\int_0^1 f(\xdel{\alpha})\,d\alpha
\qquad
\mathrm{AOPC}_{del}=\frac1K\sum_{k=1}^K\left[f(x)-f(\xdel{k})\right]
\]
\centering
\begin{tikzpicture}[x=1cm,y=1cm,>=stealth]
\tikzset{axis/.style={draw=black!70,line width=.35pt}, lab/.style={font=\scriptsize,align=center}}
\begin{scope}[shift={(0,0)}]
  \node[lab,font=\bfseries\scriptsize] at (1.55,2.15) {$\mathrm{AUC}_{del}$};
  \draw[axis] (0,0) -- (0,1.75) -- (3.1,1.75);
  \node[lab,rotate=90] at (-.35,.95) {score};
  \node[lab] at (1.55,-.28) {features deleted};
  \fill[contentblue!18] (0,0) -- (0,1.55) .. controls (.7,1.32) and (1.2,.85) .. (1.75,.55)
    .. controls (2.25,.32) and (2.7,.24) .. (3.05,.2) -- (3.05,0) -- cycle;
  \draw[contentblue,line width=1pt] (0,1.55) .. controls (.7,1.32) and (1.2,.85) .. (1.75,.55)
    .. controls (2.25,.32) and (2.7,.24) .. (3.05,.2);
  \node[lab,contentblue] at (1.65,.55) {area under\\deletion curve};
\end{scope}
\begin{scope}[shift={(4.5,0)}]
  \node[lab,font=\bfseries\scriptsize] at (1.55,2.15) {$\mathrm{AOPC}_{del}$};
  \draw[axis] (0,0) -- (0,1.75) -- (3.1,1.75);
  \node[lab,rotate=90] at (-.35,.95) {drop};
  \node[lab] at (1.55,-.28) {deletion step $k$};
  \foreach \x/\h in {.35/.25,.8/.55,1.25/.86,1.7/1.08,2.15/1.22,2.6/1.32}
    \draw[fill=contentblue!35,draw=contentblue!85] (\x,0) rectangle +(0.26,\h);
  \draw[contentblue,line width=1pt] (.25,1.0) -- (2.95,1.0);
  \node[lab,contentblue] at (2.25,1.35) {mean\\drop};
\end{scope}
\end{tikzpicture}

\vspace{.4em}
\raggedright\small For a score-preserving output and a fixed deletion rule, lower deletion AUC or higher AOPC is generally preferred. Always state direction and target output.
\end{frame}

\begin{frame}{Comprehensiveness and sufficiency}
\HierarchyHeader{\ContentIcon}{Content}{Completeness}{Is the selected evidence enough, and is removing it consequential?}
\[
\mathrm{Comp}(x,S)=f(x)-f(\xdel{S})\qquad
\mathrm{Suff}(x,S)=|f(x)-f(\xkeep{S})|
\]
\begin{itemize}
  \item Remove selected features: does the score change? Higher drop is generally preferred.
  \item Keep selected features: does the score remain? Smaller difference is generally preferred.
\end{itemize}
\centering\includegraphics[width=.66\linewidth,height=.3\textheight,keepaspectratio]{comprehensiveness_sufficiency.pdf}
\end{frame}

\begin{frame}{Infidelity}
\HierarchyHeader{\ContentIcon}{Content}{Correctness / faithfulness}{Does the attribution-predicted change match the real change?}
\[
\operatorname{Infid}(E,f,x)=\mathbb{E}_{\delta}\left[\left(\delta^\top E(f,x)-\left(f(x)-f(x-\delta)\right)\right)^2\right]
\]
\centering
\begin{tikzpicture}[x=1cm,y=1cm,>=stealth]
\tikzset{
  box/.style={draw=black!45,fill=black!3,minimum width=2.1cm,minimum height=.58cm,align=center,font=\scriptsize},
  lab/.style={font=\scriptsize,align=center},
  arr/.style={->,draw=black!65,line width=.45pt}
}
\node[box] (pert) at (0,1.1) {perturbation\\$\delta$};
\node[box,fill=contentblue!8] (pred) at (3.0,1.75) {attribution\\prediction\\$\delta^\top E(f,x)$};
\node[box,fill=contentblue!8] (actual) at (3.0,.45) {model\\change\\$f(x)-f(x-\delta)$};
\node[box,fill=contentblue!14,minimum width=2.55cm] (gap) at (6.2,1.1) {squared\\mismatch};
\draw[arr] (pert.east) -- (pred.west);
\draw[arr] (pert.east) -- (actual.west);
\draw[arr] (pred.east) -- (gap.west);
\draw[arr] (actual.east) -- (gap.west);
\node[lab,contentblue] at (6.2,.25) {small gap $\Rightarrow$ low infidelity};
\end{tikzpicture}

\vspace{.35em}
\raggedright\small
It compares attribution-predicted change with actual model change. The perturbation distribution over $\delta$ defines what changes are tested.\cite{infidelity}
\end{frame}

\begin{frame}{The perturbation problem}
\HierarchyHeader{\ContentIcon}{Content}{Measurement validity}{Does the perturbation still represent a meaningful input?}
\centering\includegraphics[width=.5\linewidth,height=.45\textheight,keepaspectratio]{ood_perturbation.pdf}
\begin{itemize}
  \item What does ``feature absence'' mean: zero, mean, median, or conditional sample?
  \item If replacement leaves the data manifold, are we measuring the explanation or model extrapolation?
\end{itemize}
\end{frame}

\begin{frame}{Baseline sensitivity is a result}
\HierarchyHeader{\ContentIcon}{Content}{Measurement validity}{How much does the score depend on the replacement convention?}
\centering\includegraphics[width=.74\linewidth,height=.55\textheight,keepaspectratio]{baseline_sensitivity.pdf}
\small Measured on one held-out instance; the curve changes when the replacement convention changes.
\end{frame}
